\documentclass[a4paper,10pt]{article}
\usepackage[utf8]{inputenc}

% ----  Useful packages % ---- 
\usepackage{amsmath}
\usepackage{graphicx}
\usepackage{amsfonts}
\usepackage{amsthm}
\usepackage{amssymb}
\usepackage{makecell}
% ----  Useful packages % ---- 

\usepackage{wrapfig}
\usepackage{caption}
\usepackage{subcaption}
\usepackage{xcolor}
\usepackage{hyperref}
\hypersetup{
    colorlinks,
    citecolor=black,
    filecolor=black,
    linkcolor=black,
    urlcolor=black
}

\graphicspath{ {./images/} }

% ---- Set page size and margins replace ------
\usepackage[letterpaper,top=2cm,bottom=2cm,left=3cm,right=3cm,marginparwidth=1.75cm]{geometry}
% ---- Set page size and margins replace ------

% ------- NOTA ------
\theoremstyle{remark}
\newtheorem{note}{Note}[subsection]
% ------- NOTA ------

% ------- OSSERVAZIONE ------
\theoremstyle{definition}
\newtheorem{observation}{Osservazione}[subsection]
% ------- OSSERVAZIONE ------

% ------- DEFINIZIONE ------
\theoremstyle{plain}
\newtheorem{definition}{Definizione}[subsection]
% ------- DEFINIZIONE ------

% ------- ESEMPIO ------
\theoremstyle{definition}
\newtheorem{example}{Esempio}[subsection]
% ------- ESEMPIO ------

% ------- DIMOSTRAZIONE ------
\theoremstyle{definition}
\newtheorem{demostration}{Dimotrazione}[subsection]
% ------- DIMOSTRAZIONE ------

% ------- TEOREMA ------
\theoremstyle{definition}
\newtheorem{theorem}{Teorema}[subsection]
% ------- TEOREMA ------

% ------- COROLLARIO ------
\theoremstyle{plain}
\newtheorem{corollaries}{Corollario}[theorem]
% ------- COROLLARIO ------

% ------- PROPOSIZIONE ------
\theoremstyle{plain}
\newtheorem{proposition}{Proposizione}[subsection]
% ------- PROPOSIZIONE ------

% ---- Footer and header ---- 
\usepackage{fancyhdr}
\pagestyle{fancy}
\fancyhf{}
\fancyhead[LE,RO]{A.A 2022-2023}
\fancyhead[RE,LO]{Analisi Matematica}
\fancyfoot[RE,LO]{\rightmark}
\fancyfoot[LE,RO]{\thepage}

\renewcommand{\headrulewidth}{.5pt}
\renewcommand{\footrulewidth}{.5pt}
\newcolumntype{P}[1]{>{\centering\arraybackslash}p{#1}}
% ---- Footer and header ---- 

% ----  Language setting ---- 
\usepackage[italian, english]{babel}
% ----  Language setting ---- 

\title{\textbf{Analisi Matematica}}
\author{Realizzato da: Giuntoni Matteo e Ghirardini FIlippo}
\date{A.A. 2022-2023}

\begin{document}
\begin{titlepage} %crea l'enviroment
\begin{figure}[t] %inserisce le figure
    \centering\includegraphics[width=0.98\textwidth]{marchio_unipi_pant541.png}
\end{figure}
\vspace{20mm}

\begin{Large}
 \begin{center}
	\textbf{Dipartimento di Informatica\\ Corso di Laurea Triennale in Informatica\\}
	\vspace{20mm}
    {\LARGE{Corso 1° anno - 6 CFU}}\\
	\vspace{10mm}
	{\huge{\bf Algebra Lineare}}\\
\end{center}
\end{Large}


\vspace{36mm}
%minipage divide la pagina in due sezioni settabili
\begin{minipage}[t]{0.47\textwidth}
	{\large{\bf Professore:}\\ \large{Prof. Tamas Szamuely}}
\end{minipage}
\hfill
\begin{minipage}[t]{0.47\textwidth}\raggedleft
	{\large{\bf Autore:}\\ \large{Matteo Giuntoni} \\ \large{Filippo Ghirardini}}
\end{minipage}

\vspace{25mm}

\hrulefill

\vspace{5mm}

\centering{\large{\bf Anno Accademico 2022/2023 }}

\end{titlepage}

\tableofcontents
\newpage
\maketitle
\begin{center}
    \vspace{-20pt}
    \rule{11cm}{.1pt} 
\end{center}

\input{introduzione}
\input{funzioni}
\input{massimi-minimi}
\input{valore-assoluto}
\input{continuità}
\input{limiti}
\input{infinitesimi}
\input{asintoti}
\input{derivate}
\input{taylor}
\input{convessità}
\input{integrali}
\input{successioni}
\input{serie}
\input{più-variabili}
\input{limiti-più-variabili}
\input{calcolo-differenziale-più-variabili}
\input{max-min-più-variabili}
% !TeX spellcheck = it_IT
\newpage
\section{Cheatsheet}

\begin{table}[h!]
	\setlength{\tabcolsep}{7pt}
	\renewcommand{\arraystretch}{2}
	\centering
	\begin{tabular}{|c|P{180px}|}
		\hline
		\textbf{Disuguaglianza triangolare} & \multirowcell{2}{ $|a + b| \leq |a| + |b|$\\$||a| + |b|| \leq |a - b| $} \\
		\hline
		\textbf{Rapporto incrementale} & $\frac{\Delta y}{\Delta x} = \frac{f(x_1)-f(x_0)}{x_1-x_0}$ \\
		\hline
		\textbf{Funzione pari e dispari} & \multirowcell{2}{$f(-x) = f(x)$ \\ $f(-x) = -f(x)$} \\
		\hline
		\textbf{Boh} & $f(x)^{g(x)} = e^{{\log{(f(x)}^ {g(x)})}} = e^{ g(x) \cdot \log{(f(x))}}$ \\
		\hline
		\textbf{Continuità in un punto} & $\lim_{x \rightarrow {x_0}^+} f(x) = \lim_{x \rightarrow {x_0}^-} f(x) = f(x_0)$ \\
		\hline
		\textbf{Derivabilità in un punto} & $\lim_{x \rightarrow {x_0}^+} f'(x) = \lim_{x \rightarrow {x_0}^-} f'(x)$\\
		\hline
		\textbf{Asintoto obliquo} & \multirowcell{2}{
			$m = \lim_{x \rightarrow \pm\infty} \frac{f(x)}{x}$ \\
			$q = \lim_{x \rightarrow \pm\infty} (f(x) -x)$
		}\\
		\hline
		
	\end{tabular}
	\caption{Formule varie}
\end{table}

\begin{table}[h!]
	\renewcommand{\arraystretch}{2.15}
	\centering
	\begin{tabular}{|c|P{210px}|}
		\hline
		\textbf{Zeri} &  \multirowcell{2}{\color{blue} $f: [a, b] \longrightarrow \mathbb{R}$ continua\\ \color{red}$f(a) \cdot f(b) < 0 \Longrightarrow \exists c \in (a, b) : f(c) = 0$ } \\
		\hline
		\textbf{Confronto} & \multirowcell{5}{\color{blue}$A \subset \mathbb{R}$, $x_0 \in Acc(x)$, $f,g: A \to \mathbb{R}$ \\ \color[HTML]{00A64F}$\lim\limits_{x\to x_0}f(x) = l_1 \wedge \lim\limits_{x\to x_0}g(x) = l_2 \wedge $ \\ \color[HTML]{00A64F}$\exists U \text{ intorno } x_0:x \in U \cap A \setminus \{x_0\} $ \\ \color{red}$f(x) \leq g(x) \Longrightarrow l_1 \leq l_2$ \\ \color{red}$f(x) \geq g(x) \Longrightarrow l_1 \geq l_2$}\\
		\hline
		\textbf{Weirstrass} & \multirowcell{7}{
			\color{blue}$a,b \in \overline{\mathbb{R}}$, $f: (a,b) \to \mathbb{R}$ continua:\\
			\color[HTML]{00A64F}$\exists \: \lim\limits_{x \to a}f(x) = l_1 \wedge \exists \lim\limits_{x \to b}f(x) = l_2$\\
			\color{red}$f$ lim. inf. $\Longleftrightarrow$ $l_1 \neq -\infty \wedge l_2 \neq -\infty$ \\
			\color{red}$f$ lim. sup. $\Longleftrightarrow$ $l_1 \neq +\infty \wedge l_2 \neq +\infty$\\
			\color{red}$f$ lim. $\Longleftrightarrow$ $l_1 \in \mathbb{R} \wedge l_2 \in \mathbb{R}$\\
			\color{red}$f$ ha min $\Longleftrightarrow \: \exists x_0 \in (a,b) : f(x_0) \leq min\{l_1, l_2\}$\\
			\color{red}$f$ ha max $\Longleftrightarrow \: \exists x_0 \in (a,b) : f(x_0) \geq max\{l_1, l_2\}$
		}\\
		\hline
		\textbf{Carabinieri} & Se due funzioni hanno lo stesso limite ed una è inferiore all'altra, se esiste una $g(x)$ in mezzo a queste due, avrà lo stesso limite\\
		\hline
		\textbf{Lagrange} & \multirowcell{4}{
			\color{blue}$f: [a,b] \to \mathbb{R}$ continua \\
			\color{blue} $f$ derivabile in $(a,b)$ \\
			\color{red} $\exists c \in (a,b): f'(c) = \frac{f(b) - f(a)}{b - a}$
		}\\
		\hline
		\textbf{De l'Hopital} & \multirowcell{3}{
			\color{blue} $\lim\limits_{x\to a^+}f(x) = \lim\limits_{x\to a^+}g(x) = 0$ oppure \\
			\color{blue} $\lim\limits_{x\to a^+}f(x) = \pm\infty$ e $\lim\limits_{x\to a^+}g(x) = \pm\infty$ \\
			\color[HTML]{00A64F} $g'(x) \neq 0$ in un introno destro di $a$ \\
			\color{red} $\lim\limits_{x\to a^+}\frac{f(x)}{g(x)} = \lim\limits_{x\to a^+}\frac{f'(x)}{g'(x)}$
		}\\
		\hline
		\textbf{o-piccolo} & \color{red}$f(x) = o(g(x)) \Longleftrightarrow \lim_{x \rightarrow x_0} \frac{f(x)}{g(x)} = 0$ \\
		\hline
		\textbf{O-grande} & \multirowcell{5}{
			\color{blue}$A \subset \mathbb{R} \wedge x_0 \in Acc(A) \wedge f,g: A \to \mathbb{R}$ \\
			\color{blue}$U$ intorno di $x_0$ \\
			\color{blue}$\exists M \in \mathbb{R}:$\\
			\color{red}$|f(x)| \geq M \cdot |g(x)| \Longrightarrow f(x) = O(g(x))$ \\
			\color[HTML]{00A64F}$\forall x \in U \cap A \setminus \{x_0\}$
		}\\
		\hline
		\textbf{Infinitesima e parte principale} & \multirowcell{4}{
			\color{blue}$A \subset \mathbb{R} \wedge x_0 \in Acc(A)$\\
			\color{blue}$f,g: A \to \mathbb{R}$ infinitesime per $x\to x_0$ \\
			\color[HTML]{00A64F}$\exists L, \alpha \in \mathbb{R}$ con $L \neq 0$ t.c. $f(x) = L \cdot (g(x))^\alpha + o((g(x))^\alpha)$ \\
			\color{red} $f$ infinitesimo di ordine $\alpha$ rispetto a $g$ \\
			\color{red} $L(g(x))^\alpha$ parte principale
		}\\
	\hline
	\end{tabular}
	\caption{Teoremi e definizioni}
\end{table}

\begin{table}[h!]
	\setlength{\tabcolsep}{7pt}
	\renewcommand{\arraystretch}{1.5}
	\centering
	\begin{tabular}{|c c c|}
		\hline
		$[1]$ $(+\infty) + (-\infty)$ & $[2]$ $(-\infty) + (+\infty)$ & $[3]$ $0 \cdot (\pm \infty)$ \\
		$[4]$ $(\pm \infty)^0$ & $[5]$ $(0^+)^0$ & $[6]$ $(1)^{\pm \infty}$\\ 
		\hline
	\end{tabular}
	\caption{Forme indeterminate}
\end{table}

\begin{table}[h!]
	\setlength{\tabcolsep}{7pt}
	\renewcommand{\arraystretch}{1.5}
	\centering
	\begin{tabular}{|c|c|}
		\hline
		\textbf{SE} & \textbf{ALLORA} \\
		\hline
		$\lim\limits_{x\to x_0}f(x) = 0^+$ & $\lim\limits_{x\to x_0}\frac{1}{f(x)} = +\infty$ \\
		\hline
		$\lim\limits_{x\to x_0}f(x) = 0^-$ & $\lim\limits_{x\to x_0}\frac{1}{f(x)} = -\infty$ \\
		\hline
		$\lim\limits_{x\to x_0}f(x) = +\infty$ & $\lim\limits_{x\to x_0}\frac{1}{f(x)} = 0^+$ \\
		\hline
		$\lim\limits_{x\to x_0}f(x) = -\infty$ & $\lim\limits_{x\to x_0}\frac{1}{f(x)} = 0^-$ \\
		\hline
		$\lim\limits_{x\to x_0}f(x) = l$ con $l \neq 0, \pm\infty$ & $\lim\limits_{x\to x_0}\frac{1}{f(x)} = \frac{1}{l}$ \\
		\hline
	\end{tabular}
	\caption{Limiti ovvi}
\end{table}

\begin{table}[h!]
	\setlength{\tabcolsep}{7pt}
	\renewcommand{\arraystretch}{1.5}
	\centering
	\begin{tabular}{|c|c c|}
		\hline
		$\mathbf{x^n}$ & $\lim\limits_{x\to +\infty}x^n = +\infty$ & $\lim\limits_{x\to +\infty}\frac{1}{x^n} = \frac{1}{+\infty} = 0$ \\
		\hline
		$\mathbf{e^x }$& $\lim\limits_{x\to +\infty}e^x = +\infty$ & $\lim\limits_{x\to -\infty}e^x = 0^+$ \\
		\hline
		$\mathbf{\log x}$ & $\lim\limits_{x\to 0^+}\log(x) = -\infty$ & $\lim\limits_{x\to +\infty}\log(x) = +\infty$ \\
		\hline
		$\mathbf{a^x}, a\geq 1$ &  $\lim\limits_{x\to +\infty}a^x = +\infty$ & $\lim\limits_{x\to -\infty}a^x = 0^+$ \\
		\hline
		$a = 1$ & $\lim\limits_{x\to +\infty}a^x = 1$ & $\lim\limits_{x\to -\infty}a^x = 1$  \\
		\hline
		$0 < a < 1$ & $\lim\limits_{x\to +\infty}a^x = 0^+$ & $\lim\limits_{x\to -\infty}a^x = +\infty$\\
		\hline
	\end{tabular}
	\vspace{-5pt}
	\caption{Limiti fondamentali}
\end{table}

Limite di un polinomio che tende ad infinito: raccoglimento\\
Limite di rapporto di polinomi che tende ad infinito: raccoglimento

\begin{table}[h!]
	\centering
	\setlength{\tabcolsep}{10pt}
	\renewcommand{\arraystretch}{2.5}
	\begin{tabular}{|c|c|}
		\hline
		$\lim\limits_{x\to 0}\frac{\sin(x)}{x} = 1$ & $\lim\limits_{x\to 0} \frac{1-\cos(x)}{x^2} = \frac{1}{2}$ \\
		\hline
		$\lim\limits_{x\to 0}\frac{e^x-1}{x} = 1$ & $\lim\limits_{x\to 0}\frac{\log(1+x)}{x} = 1$\\
		\hline
	\end{tabular}
	\caption{Limiti notevoli}
\end{table}

\begin{table}[h!]
\centering
\setlength{\tabcolsep}{10pt}
\renewcommand{\arraystretch}{2.5}
\begin{tabular}{|c|}
	\hline
	$\lim\limits_{x \to +\infty}\frac{\log(x)}{x} = \lim\limits_{y\to +\infty}\frac{y}{e^y} = 0$ \\
	\hline
	$\lim\limits_{x\to +\infty}\frac{(\log(x))^\beta}{x^\alpha} = \lim\limits_{y \to +\infty}\frac{y^\beta}{(e^y)^\alpha} = \lim\limits_{y \to +\infty}\frac{y^\beta}{e^{y\cdot\alpha}} = 0$  \\
	\hline
	$\lim\limits_{x\to 0^+}x\log(x) = \lim\limits_{y\to -\infty}e^y \cdot y = \lim\limits_{z\to +\infty}e^{-z} \cdot (-z) = \frac{-z}{e^z} = 0$ \\
	\hline
	$\lim\limits_{x\to 0^+}x^\alpha \cdot \log(x) = \lim\limits_{y\to 0^+}y \cdot \log(y^{\frac{1}{\alpha}}) = \lim\limits_{y\to 0^+}\frac{y}{\alpha} \cdot \log(y) = \frac{1}{\alpha}\lim\limits_{y\to 0^+} y \cdot \log(y) = 0$ \\
	\hline
\end{tabular}
\caption{Limiti con sostituzione}
\end{table}

\begin{table}[h!]
	\centering
	\setlength{\tabcolsep}{10pt}
	\renewcommand{\arraystretch}{2.5}
	\begin{tabular}{|c|c|}
		\hline
		\textbf{Prodotto} & $f(x) \cdot o(g(x)) = o(f(x) \cdot g(x))$\\
		\hline
		\textbf{Prodotto costante} & $k \in \mathbb{R}, k \neq 0 \Longrightarrow o(k \cdot g(x)) = o(g(x))$\\
		\hline
		\textbf{Somma} & $o(g) + o(g) = o(g)$ \\
		\hline
		\textbf{Prodotto 0} & $\lim\limits_{x\to x_0}f(x) = 0 \Longrightarrow f(x) \cdot g(x) = o(g(x))$ \\
		\hline
		\textbf{Somma 0} & $\lim\limits_{x\to x_0}f(x) = 0 \Longrightarrow o(g) + o(f \cdot g) = o(g)$ \\
		\hline
		\textbf{o di o} & $o(o(g)) = o(g)$ \\
		\hline
		\textbf{o di somma} & $o(f + g) = o(f) + o(g)$ \\
		\hline
		\textbf{o di prodotto} & $o(g) \cdot o(f) = o(f \cdot g)$ \\
		\hline
	\end{tabular}
	\caption{Proprietà o-piccolo}
\end{table}
\begin{table}[h!]
	\setlength{\tabcolsep}{5pt}
	\renewcommand{\arraystretch}{2.2}
	\centering
	\begin{tabular}{|c|c|}
		\hline
		$e^x$ & $1 + x + \frac{x^2}{2!} + \frac{x^3}{3!} + \frac{x^4}{4!} + ... + \frac{x^n}{n!} + o(x^n)$  \\
		$\log(1+x)$ & $x - \frac{x^2}{2} + \frac{x^3}{3} - \frac{x^4}{4} + \frac{x^5}{5} + ... + (-1)^{n-1}\frac{x^n}{n} + o(x^n)$ \\
		$\sin(x)$ & $x - \frac{x^3}{3!} + \frac{x^5}{5!} - \frac{x^7}{7!} + ... + (-1)^n \frac{x^{2x+1}}{(2n+1)!} + o(x^{2n+2})$ \\
		$\cos(x)$ & $1 - \frac{x^2}{2!} + \frac{x^4}{4!} - \frac{x^6}{6!} + ... + (-1)^n\frac{x^2n}{(2n)!} + o(x^{2n+1})$ \\
		$\tan(x)$ & $x + \frac{x^3}{3} + \frac{2}{15}x^5 + o(x^6)$\\
		$\arctan(x)$ & $x - \frac{x^3}{3} + \frac{x^5}{5} - \frac{x^7}{7} + ... + (-1)^n\frac{x^{2x+1}}{(2n + 1)} + o(x^{2n+2})$\\
		$\arcsin{x}$ & $x + \frac{x^3}{6} + \frac{3}{40}x^5 + o(x^6)$\\
		$\sqrt{1+x}$ & $1 + \frac{1}{2}x - \frac{1}{8}x^2 + \frac{1}{16}x^3 + o(x^3)$\\
		$(1+x)^{\alpha}$ & $1 + \alpha x + \frac{\alpha(\alpha - 1)}{2}x^2 + \frac{\alpha(\alpha - 1)(\alpha - 2)}{6}x^3 + o(x^3)$\\
		\hline
	\end{tabular}
	\caption{Formule di taylor}
\end{table}

\begin{table}[h!]
	\setlength{\tabcolsep}{5pt}
	\renewcommand{\arraystretch}{2.2}
	\centering
	\begin{tabular}{|c|c|}
		\hline
		$\lim_{x \rightarrow +\infty} \frac{a^x}{x^\alpha} = + \infty$ se $a>1$ \\
		\hline
		$\lim_{x \rightarrow +\infty} \frac{a^x}{x^\alpha} = 0^+$ se $0<a<1$ \\
		\hline
		$\lim_{x \rightarrow +\infty} \frac{{\log{x}}^\beta}{x^\alpha} = 0$ \\
		\hline
		$\lim_{x \rightarrow 0^+} x^\alpha \cdot \log{x} = 0$ \\
		\hline
	\end{tabular}
	\caption{Confronto tra infiniti e infinitesimi}
\end{table}

Derivate fondamentali
Inclusa composta e inversa
\end{document}
